\begin{abstract}
Single-cell transcriptomics data are sparse and noisy; among approaches to address this, cell aggregation methods group similar cells into representative profiles termed metacells. Yet current methods define similarity from transcriptomics alone and operate within individual batches: in multi-batch settings, batch effects limit cross-batch aggregation even when biological states are shared; moreover, niche-specific transcriptional programs encoded in spatial organization are invisible to transcriptomic similarity, and methods that optimize purely on affinity structure cannot incorporate spatial context without sacrificing cell-type identity. We introduce scProto, a framework that learns metacells by combining a reconstruction objective with a community-preserving objective that aligns prototype assignments with a cell--cell affinity graph encoding biologically meaningful similarity (e.g., density-corrected transcriptomic or spatial). Experiments on single-cell and spatial transcriptomics datasets show that scProto improves cross-batch aggregation while preserving community structure and cell-type purity, better captures rare cell populations, and---by incorporating spatial affinity---identifies niche-correlated transcriptional programs while preserving cell-type identity.
\end{abstract}
